\documentclass[a4paper,12pt]{article}
\usepackage{fontspec}
\usepackage{polyglossia}
\setotherlanguage{english}
\setmainfont{LiberationSans} % or another font that supports Greek
\usepackage{amsmath, amssymb}
\usepackage[most]{tcolorbox}
\usepackage{xcolor}
\usepackage{colortbl}
\usepackage{enumitem}
\usepackage{multirow}
\usepackage[a4paper, margin=1.8cm]{geometry}
\usepackage{fancyhdr}
\pagestyle{fancy}
\renewcommand{\headrulewidth}{0pt}
\fancyhf{} 
\fancyfoot[L]{\footnotesize Theolaos}
\fancyfoot[R]{\small \thepage}

\usetikzlibrary{decorations.pathreplacing}
\usetikzlibrary{arrows.meta}

\setlist[itemize]{left=1.5em, label=\textbullet}
\setlength{\parskip}{0pt}
\setlength{\parindent}{0pt}

% ΕΡΓΑΣΙΑ_1_Ερ-ΜΚΕΡ-ΣΑΒΒΙΔΗΣ_ΘΕΟΦΙΛΟΣ_ΝΙΚΟΛΑΟΣ

\begin{document}

\begin{center}
    \LARGE \textbf{Cheat Sheet}\\[1ex]
    \large Ολοκληρωμάτων
\end{center}

\section*{Μέθοδοι Ολοκλήρωσης}

\subsection*{Κατά παράγοντες:}
\[
\int f(x)g(x)\,dx = \int F'(x)g(x)\,dx = F(x)g(x) - \int F(x)g'(x)\,dx
\]

\subsection*{Με αντικατάσταση:}
Έστω $f(x) = u$, τότε:
\[
\frac{du}{dx} = u' \iff dx = \frac{du}{u'}
\]
και απλά αντικαθιστώ στη συνέχεια στο ολοκλήρωμα.

Κάποιες φορές στα κλάσματα μπορώ να αντικαταστήσω με σκοπό να βγάλω τη μεταβλητή στον αριθμητή ή και στον παρονομαστή.

\paragraph{Παράδειγμα:}
\[
\int \frac{x}{\sqrt{1-x^2}}\,dx
\]
Έστω $u = 1-x^2$, οπότε $du = -2x\,dx \iff dx = -\frac{du}{2x}$.

Άρα:
\[
\int \frac{x}{\sqrt{u}} \cdot \frac{du}{-2x} = \int \frac{1}{\sqrt{u}} \cdot \frac{du}{-2} = -\int \frac{du}{2\sqrt{u}}
\]
Και φυσικά η λύση θα είναι:
\[
-\sqrt{u} + C
\]

\subsection*{Τριγωνομετρικά Ολοκληρώματα}

\subsubsection*{Περίπτωση μονών δυνάμεων:}
Σε περίπτωση που έχουμε μονές δυνάμεις, π.χ. $\int f(\sin x, \cos x)\,dx$, τότε μπορούμε να εφαρμό\-σουμε την παρακάτω αντικατάσταση:
\[
u = \tan\frac{x}{2}
\]
\[
dx = \frac{2\,du}{1 + u^2}, \quad \sin x = \frac{2u}{1+u^2}, \quad \cos x = \frac{1-u^2}{1+u^2}, \quad \tan x = \frac{2u}{1-u^2}
\]

\subsubsection*{Περίπτωση δυνάμεων εις την δευτέρα:}
Σε περίπτωση που έχεις δύναμη εις την δύο, π.χ. $\int f(\sin^2 x, \cos^2 x)\,dx$, τότε κάνεις την παρακά\-τω αντικατάσταση:
\[
u = \tan x
\]
\[
dx = \frac{du}{1+u^2}, \quad \sin^2 x = \frac{u^2}{1+u^2}, \quad \cos^2 x = \frac{1}{1+u^2}
\]

\subsection*{Τριγωνομετρικές Ταυτότητες}
\[
\sin^2 x = \frac{1}{2}(1-\cos 2x)
\]
\[
\cos^2 x = \frac{1}{2}(1+\cos 2x)
\]
\[
\sin x \cos y = \frac{1}{2}(\sin(x-y) + \sin(x+y))
\]
\[
\cos x \cos y = \frac{1}{2}(\cos(x-y) + \cos(x+y))
\]
\[
\sin x \sin y = \frac{1}{2}(\cos(x-y) - \cos(x+y))
\]
\[
\cos(x+y) = \cos x \cos y - \sin x \sin y
\]
\[
\sin(x+y) = \sin x \cos y + \cos x \sin y
\]

\subsection*{Βασικά Αόριστα Ολοκληρώματα}
Κάποια εξτρά χρήσιμα βασικά αόριστα ολοκληρώματα, πέρα από τα γνωστά, είναι:
\[
\int a^x\,dx = \frac{a^x}{\ln a} + C
\]
\[
\int \frac{1}{\cos^2 x}\,dx = \tan x + C
\]
\[
\int \frac{1}{\sqrt{1-x^2}}\,dx = \arcsin x + C
\]
\[
\int \frac{1}{1+x^2}\,dx = \arctan x + C
\]

\end{document}
